\section*{ABSTRACT:}
\noindent

Modern digital memory systems---including DRAM, NAND flash, and on-chip SRAM---are increasingly susceptible to bit-flip errors as device geometries shrink and supply voltages are scaled for energy efficiency. These errors corrupt stored data and, in safety-critical applications such as autonomous vehicle inference or aerospace computing, can cause silent failures with catastrophic consequences. The dominant error-correcting code (ECC) in deployed systems, Bose-Chaudhuri-Hocquenghem (BCH), requires Galois field arithmetic for decoding and imposes metadata overhead that grows steeply with the number of correctable errors. At high bit error rates ($\geq$1\%), BCH overhead becomes impractical, and its algebraic structure provides no native defense against burst errors---contiguous runs of flipped bits common in flash storage and DRAM row-hammer events.

This disclosure presents [ACRONYM], a novel approximate error correction scheme that corrects bit-flip errors in data blocks without Galois field arithmetic, with metadata overhead that \emph{decreases} as block size increases. [ACRONYM] encodes a data block by first applying a keyed, invertible Feistel permutation to its bit indices, scattering any burst errors uniformly across a two-dimensional (2D) grid of $N \times N$ cells (where $N \approx \sqrt{L}$ for a block of $L$ bits). Cyclic redundancy check (CRC) digests are then computed over non-overlapping groups of rows and columns of the grid and stored as compact error-correction metadata alongside the data.

During decoding, [ACRONYM] recomputes row and column digests over a received block, identifies mismatched hash nodes, and constructs a constraint graph in which nodes represent digest computations and edges connect nodes sharing bit coverage. A greedy, constraint-guided bit-flip search then iteratively corrects errors by enumerating flip combinations within a pruned candidate set---excluding bits already ``pinned'' as correct by neighboring constraints---and committing the combination that maximizes total matched nodes. This matched-neighbor pinning heuristic is the key algorithmic innovation that makes exhaustive combinatorial search tractable even at 5\% bit error rate.

Empirical results on 4,096-bit blocks demonstrate that [ACRONYM] with CRC-32 hashes (100\% overhead) fully corrects 200 or more bit-flips with $\geq$95\% success rate, while BCH codes require approximately 174\% overhead to achieve equivalent correction capability. Overhead scales as $O(1/\sqrt{L})$, dropping to $\sim$50\% at 16,384-bit blocks with the same hash configuration. Random and burst errors are corrected with statistically identical success rates and search effort, demonstrating complete burst resilience. [ACRONYM] provides a practical, hardware-friendly, and tunable ECC solution particularly well suited to protecting deep neural network (DNN) weight memory in AI accelerators, NAND flash storage, and high-density DRAM systems.
